\documentclass[a4paper, 12pt]{mksheetphhd}

\title{Listen über Listen}
\author{Marvin Ködding}
\usepackage{minted}
\usepackage{multicol}

\everymath{\displaystyle}

\begin{document}
	\printtitle
	
	\vspace*{1cm}
	
	\begin{aufgabe}
		Man kann Bedingungen, die man testen will, miteinander mittels Junktoren verknüpfen. Diese sind \verb*|not|, \verb*|and| und \verb*|or|. Im folgenden Block wird der \verb*|print|-Befehl nicht ausgeführt, wenn der Wert von $x$ zwischen $10$ und $20$ liegt.
		
		\begin{minted}{python}
if x < 10 or x > 20:
  print(x)
		\end{minted}
		
		Der Block
		
		\begin{minted}{python}
if not (x >= 10 and x <= 20):
  print(x)
		\end{minted}
		
		bewirkt dabei den gleichen Effekt. Probiert es selbst aus. Verknüpft mehrere Aussagen mittels Junktoren.
	\end{aufgabe}
	
	\begin{comment}
		Berechnet das Produkt $1 \cdot 3 \cdot 5 \cdot \dots \cdot 33$ \dots
		\begin{enumerate}
			\item mit einer while-Schleife und
			\item mit einer for-Schleife
		\end{enumerate}
	\end{comment}
	
	\begin{aufgabe}
		Gebt die Summe der ersten $100$ natürlichen Zahlen an, die durch $7$ teilbar sind. Verwendet dabei
		\begin{enumerate}
			\item eine while-Schleife,
			\item eine for-Schleife,
			\item List Comprehensions.
		\end{enumerate}
		Findet andere Beispiele, bei denen ein Schleifentyp besser ist als ein anderer.
	\end{aufgabe}
	
	\begin{aufgabe}
		Ihr bekommt eine Zeichenkette übergeben. Verwendet eine List Comprehension, um zu zählen, wie viele Vokale diese Zeichenkette enthält.
	\end{aufgabe}
	
	\begin{aufgabe}
		Ihr bekommt eine Liste an Wörter übergeben. Gebt eine Liste zurück, in der die Längen der Wörter stehen. Bspw. soll zur Eingabe \mintinline{python}|["Banane", "Apfel"]| die Ausgabe \mintinline{python}|[6, 5]| zurückgegeben werden. Verwendet eine for-Schleife und List Comprehensions.
	\end{aufgabe}
	
	\textit{Bearbeitet die folgenden Aufgaben ausführlich. Diese brauchen wir in der Plenumsveranstaltung zur Diskussion.}
	
	\begin{aufgabe}
		Schreibt eine Funktion \mintinline{python}|fib(n)|, die die $n$-te Fibonaccizahl zurückgibt. Erstellt eine Liste der ersten $10$, $20$ und $50$ Fibonaccizahlen.
		
		\small Hinweis: Die Folge der Fibonaccizahlen wird rekursiv definiert als $f_0 = f_1 = 1$ und $f_n = f_{n-2} + f_{n-1}$. 
	\end{aufgabe}
	
	\begin{aufgabe}
		Überlegt euch eine Alternative, um die vorherige Aufgabe zu lösen. Verwendet \textit{nicht} die explizite Darstellung!
	\end{aufgabe}
	
	\printlicense
	
	\printsocials
	
\end{document}